\section{Resource Demo and Case Study}

The demo evaluates the artifact interface rather than ranking tokenizer
systems. It combines two public SID export paths with controlled rows that
calibrate the diagnostics. GRID supplies the open RQ-KMeans-style
path~\cite{ju2025grid}; ReSID supplies the GAOQ line and processed Musical
artifacts~\cite{liang2026resid}. Sanity tokenizers and stressors are included
only when they clarify what a diagnostic can and cannot identify.

\paragraph{Case-study protocol.}
The paper-facing case study uses the same 23,742 Musical items for both rows.
The GRID row is a controlled feature-text export: ReSID processed feature text
is embedded and passed through the official GRID-style residual k-means module.
The ReSID row is a bounded FAMAE-to-GAOQ export on the same item universe. This
protocol is intentionally modest. It gives reviewers a same-item diagnostic
contrast while avoiding the stronger, unsupported claim that the paper
reproduces every training choice of TIGER, GRID, or the full ReSID pipeline.

\begin{table}[H]
\caption{Same-item Musical artifact profiles on 23,742 items. Higher unique
full codes, D3 L1 co-occurrence prefix recall, and D4 tail unique-SID ratio are
better; lower D2 full-collision rate is better. D5a prefixes are level-wise
prefix counts, not latency. This is not a downstream leaderboard: GRID
feature-text uses ReSID processed feature text, and ReSID is a bounded GAOQ
export.}
\label{tab:musical-diagnostic}
\small
\setlength{\tabcolsep}{2.0pt}
\begin{tabular}{@{}lrrrrl@{}}
\toprule
System & Unique & D2 full & D3 L1 & D4 tail & D5a prefixes \\
\midrule
GRID feature-text & 3,749 & 0.9769 & 0.0552 & 0.3695 & 64/3440/3749 \\
ReSID bounded & 23,742 & 0.0000 & 0.1535 & 1.0000 & 32/1280/23742 \\
\bottomrule
\end{tabular}
\end{table}

Table~\ref{tab:musical-diagnostic} compares two exports on the same 23,742
Musical items. The GRID row is intentionally labeled as feature-text: it uses
processed feature text from the ReSID artifact path and the official
MiniBatchKMeans-style residual quantization module, not a faithful raw-text
TIGER reproduction. Under that controlled setup, the row exposes high collision
pressure and limited tail capacity. The ReSID/GAOQ row exports one full code
per item in this bounded run, yielding zero full-code collisions and higher
D3-L1 collaborative prefix recovery. The D5a prefix column reports level-wise
prefix counts, a structural proxy for trie expansion rather than real serving
latency.
The same GRID Musical feature-text export was rerun locally with seeds 43 and
44; across seeds 42--44, duplicate-\sid rate remains 0.8327--0.8421 and
full-collision rate remains 0.9751--0.9769, so the collision pressure in
Table~\ref{tab:musical-diagnostic} is not a single-seed artifact.

\begin{table}[H]
\caption{Controlled stressor signals. These rows are not named-tokenizer
evidence; they test whether D2, D4, and D5a react to specific artifact
mechanisms.}
\label{tab:controller-stressors}
\small
\setlength{\tabcolsep}{3pt}
\begin{tabular}{@{}l l r r@{}}
\toprule
Stressor & Metric & Signal 1 & Signal 2 \\
\midrule
Qualified collision & co-occurrence lift & 3.86x & 1.19x \\
Capacity budget & unique ratio & 1.000 head & 0.028 tail \\
Variable depth & prefix-cost view & active & max-depth \\
\bottomrule
\end{tabular}
\end{table}

\paragraph{Diagnostic findings.}
Tables~\ref{tab:musical-diagnostic} and~\ref{tab:controller-stressors} support
three artifact-level findings. First, collision pressure is stable in the
controlled GRID row: most items share a complete \sid{}, whereas the same
interface ingests a collision-free ReSID mapping. The qualified-collision
stressor then shows why raw collision volume is insufficient: collisions that
overlap user co-occurrence have a 3.86x lift, while synthetic hash collisions
have 1.19x. Second, collision-free addresses and collaborative-prefix alignment
are different objectives. ReSID has no full collisions and higher D3-L1 than
the GRID row, while a hand-designed category-prefix control in the repository
recovers more level-1 co-occurrence neighbors without being a learned
recommender tokenizer. Third, D4 and D5a separate capacity allocation from
prefix structure: in the capacity stressor, head uniqueness reaches 1.000
while tail uniqueness remains 0.028, and variable-depth stressors show that
active prefix cost need not match a fixed maximum-depth view.

\paragraph{Artifact-repository evidence.}
The remaining repository tables extend the same checks without changing the
paper claims. Sanity rows calibrate D1--D5a, GRID scale rows test larger item
counts, DACT rows exercise optional D6 churn, and MovieLens rows check
non-Amazon schema portability. These tables support artifact usability;
Table~\ref{tab:musical-diagnostic} remains the only method-facing case study in
the main paper.
